\section{Related Work}

\paragraph{TTA benchmarks and methods.} TTA evaluation is dominated by vision corruption benchmarks: ImageNet-C~\cite{hendrycks2019benchmarking}, CIFAR-C, WILDS~\cite{koh2021wilds}, BREEDS~\cite{santurkar2020breeds}, and Shifts~\cite{malinin2022shifts}. Representative methods include TENT~\cite{wang2021tent}, SHOT~\cite{liang2020shot}, T3A~\cite{iwasawa2021t3a}, CoTTA~\cite{wang2022cotta}, SAR~\cite{niu2023sar}, and parameter-free LAME~\cite{boudiaf2022lame}; an expanded method inventory (TTT/TTT++/NOTE/RoTTA/EATA/AdaContrast) is deferred to Appendix~\ref{app:tta_methods}. Recent work questions TTA's generality: Zhao et al.~\cite{zhao2023pitfalls} catalog pitfalls, Press et al.~\cite{press2024rdumb} show a ``dumb'' baseline competes with sophisticated methods, and Su et al.~\cite{su2024tribe} address real-world instability. All assume covariate shift~\cite{shimodaira2000covariate,quinonero2009dataset_shift}, which we argue WiFi room-to-room shift violates.

\paragraph{WiFi CSI and domain adaptation.} Standard benchmarks---SignFi~\cite{ma2018signfi}, UT-HAR~\cite{yousefi2017uthar}, Widar3.0~\cite{zheng2019widar}, SenseFi~\cite{yang2023sensefi}---and supervised or source-target domain-adaptation methods---EI~\cite{jiang2018ei}, RF-Net~\cite{ding2020rfnet}, AirFi~\cite{jiang2022airfi}, AutoFi~\cite{yang2024autofi}, DANN~\cite{ganin2016dann}, CDAN~\cite{long2018cdan}---address cross-environment robustness, but none evaluates label-free deployment-time adaptation. Calibration remains a known source of overconfidence~\cite{guo2017calibration,muller2019labelsmoothing}.

\paragraph{Physics-informed wireless.} PINNs~\cite{raissi2019pinns,karniadakis2021physics} embed PDEs as losses and can exhibit gradient pathologies~\cite{wang2021gradflow,krishnapriyan2021pinn_failures}. Recent work applies physics priors to channel estimation~\cite{javid2025nips} and radio mapping~\cite{boeck2025icml}. We evaluate whether an OFDM-derived residual prior can serve as a TTA signal.
